\section{Vista de Decisiones de Arquitectura}\label{sec:vista-de-decisiones-de-arquitectura}

La Vista de Decisiones de Arquitectura presenta y describe las principales decisiones de arquitectura tomadas \cite{VanHeesch2012}\cite{Kruchten2009}. Para cada decisión se especifican los campos de la plantilla: identificador, nombre, categorías, problema, alternativas, decisión, justificación y decisiones relacionadas.


\begin{table}[H]
    \small
    \centering
    \caption{DA01: Arquitectura en capas del componente central}
    \label{tab:da01}
    \begin{tabular}{|p{4.1cm}|p{11.7cm}|}
        \hline
            \textbf{Identificador} & DA01 \\
        \hline
            \textbf{Nombre} & Arquitectura en capas del componente central \\
        \hline
            \textbf{Categorías} & Componente central; mantenibilidad. \\
        \hline
            \textbf{Problema} & El componente central concentra la presentación para cuatro tipos de usuarios y la API del componente móvil, la lógica de negocio y de fiscalización, la persistencia y la integración con varios sistemas externos y nodos periféricos. \\
        \hline
            \textbf{Alternativas} & \pendiente{completar las alternativas consideradas.} \\
        \hline
            \textbf{Decisión} & Organizar el componente central en cuatro capas: presentación (API REST, frontoffice y backoffice), servicios y procesamiento (Servicios de Negocio, Generación de Reportes, Ingesta de Eventos y Motor de Fiscalización), persistencia, y adaptadores e infraestructura, que aíslan cada integración externa (Figura \ref{fig:vista_logica}). \\
        \hline
            \textbf{Justificación} & \pendiente{completar la justificación.} \\
        \hline
            \textbf{Decisiones Relacionadas} & DA02, DA05 \\
        \hline
    \end{tabular}
\end{table}

\begin{table}[H]
    \small
    \centering
    \caption{DA02: Broker de mensajería en su propio contenedor}
    \label{tab:da02}
    \begin{tabular}{|p{4.1cm}|p{11.7cm}|}
        \hline
            \textbf{Identificador} & DA02 \\
        \hline
            \textbf{Nombre} & Broker de mensajería en su propio contenedor \\
        \hline
            \textbf{Categorías} & Mensajería; integración con tracking; R4, R8; AC07, AC08, AC15, AC17. \\
        \hline
            \textbf{Problema} & La recepción de los eventos de tracking debe estar desacoplada de su procesamiento mediante un mecanismo de mensajería asincrónica (R4). La letra pide justificar la tecnología elegida, la estrategia de particionamiento o distribución de los mensajes y la política de entrega. \\
        \hline
            \textbf{Alternativas} & \pendiente{completar las alternativas consideradas.} \\
        \hline
            \textbf{Decisión} & Usar Apache Artemis (Jakarta Messaging) como broker en su propio contenedor, separado de WildFly y con los mensajes en un volumen. Los nodos de tracking publican en la cola \texttt{queue\_tracking} conectándose directo al broker con TLS (R8), cada uno con un usuario que solo puede publicar en esa cola. Todas las instancias del componente central consumen la misma cola, mediante el componente de recepción (Sección \ref{sec:refinamiento-ingesta}), y compiten por sus mensajes; el orden de la ruta se reconstruye por el momento de generación de cada posición (AC16) y los duplicados se descartan en la ingesta (DA07). La entrega es al menos una vez: un mensaje se borra de la cola cuando se confirma la transacción que lo procesó, y si el procesamiento falla, el broker lo vuelve a entregar. Si el procesamiento falla, la reentrega espera cada vez más (1~s, 2~s, 4~s\ldots{} hasta un minuto) y se hacen hasta 20 entregas: el componente central puede estar unos 14 minutos sin poder procesar sin perder posiciones. Recién entonces el mensaje pasa a la cola de mensajes muertos, que el tablero de la Plataforma de Observabilidad muestra y desde la que se reprocesa con la consola del broker; reprocesar es seguro porque la ingesta descarta los duplicados (DA07). \\
        \hline
            \textbf{Justificación} & \pendiente{completar la justificación.} \\
        \hline
            \textbf{Decisiones Relacionadas} & DA01, DA07 \\
        \hline
    \end{tabular}
\end{table}

\begin{table}[H]
    \small
    \centering
    \caption{DA03: Sesión sin estado en el servidor}
    \label{tab:da03}
    \begin{tabular}{|p{4.1cm}|p{11.7cm}|}
        \hline
            \textbf{Identificador} & DA03 \\
        \hline
            \textbf{Nombre} & Sesión sin estado en el servidor \\
        \hline
            \textbf{Categorías} & Seguridad; escalabilidad; AC07; CU-01. \\
        \hline
            \textbf{Problema} & El componente central debe escalar horizontalmente sin guardar estado local en las instancias (AC07). El ingreso con Usuario gub.uy necesita conservar \texttt{state}, \texttt{nonce} y \texttt{code\_challenge} entre el inicio del ingreso y el \textit{callback}. \\
        \hline
            \textbf{Alternativas} & Mantener ese estado en la \texttt{HttpSession} de la instancia y firmar el JWT de sesión con código propio, que fue la implementación inicial. \\
        \hline
            \textbf{Decisión} & La sesión es un JWT firmado (HS256) con la biblioteca JJWT, que viaja en una cookie \textit{HttpOnly} en la web y como token \textit{Bearer} en el móvil. El estado del ingreso con gub.uy viaja en una cookie firmada en lugar de quedar en la \texttt{HttpSession}. El secreto se configura por entorno y debe tener al menos 32 bytes. El JWT se emite también tras el ingreso al backoffice, que valida las credenciales contra el directorio LDAP. \\
        \hline
            \textbf{Justificación} & Hace la autenticación independiente de la instancia y saca la criptografía propia del código: el \textit{callback} de gub.uy puede llegar a cualquier instancia del componente central. \\
        \hline
            \textbf{Decisiones Relacionadas} & DA04, DA09 \\
        \hline
    \end{tabular}
\end{table}

\begin{table}[H]
    \small
    \centering
    \caption{DA04: Ingreso desde el móvil con \textit{authorization code} y PKCE}
    \label{tab:da04}
    \begin{tabular}{|p{4.1cm}|p{11.7cm}|}
        \hline
            \textbf{Identificador} & DA04 \\
        \hline
            \textbf{Nombre} & Ingreso desde el móvil con \textit{authorization code} y PKCE \\
        \hline
            \textbf{Categorías} & Seguridad; componente móvil; CU-01. \\
        \hline
            \textbf{Problema} & La aplicación móvil debe autenticar al chofer mediante Usuario gub.uy (CU-01) y obtener un token para llamar a la API REST (R2). \\
        \hline
            \textbf{Alternativas} & \pendiente{completar las alternativas consideradas.} \\
        \hline
            \textbf{Decisión} & La aplicación abre el ingreso en el navegador con un \texttt{code\_challenge}, recibe un código de un solo uso por el esquema \texttt{cargauy://ingreso} y lo canjea, junto con el \texttt{code\_verifier}, por el token de sesión. En Android, el ingreso se abre con Custom Tabs. \\
        \hline
            \textbf{Justificación} & \pendiente{completar la justificación.} \\
        \hline
            \textbf{Decisiones Relacionadas} & DA03 \\
        \hline
    \end{tabular}
\end{table}

\begin{table}[H]
    \small
    \centering
    \caption{DA05: Tecnologías de presentación y ubicación del frontoffice}
    \label{tab:da05}
    \begin{tabular}{|p{4.1cm}|p{11.7cm}|}
        \hline
            \textbf{Identificador} & DA05 \\
        \hline
            \textbf{Nombre} & Tecnologías de presentación y ubicación del frontoffice \\
        \hline
            \textbf{Categorías} & Presentación; distribución; R14. \\
        \hline
            \textbf{Problema} & La plataforma tiene un frontoffice para ciudadanos, otro para funcionarios y un backoffice para autoridades y administradores, todos web. \\
        \hline
            \textbf{Alternativas} & \pendiente{completar las alternativas consideradas.} \\
        \hline
            \textbf{Decisión} & Backoffice en JSF dentro del componente central; frontoffice en React, servido como sitio estático desde Vercel (Figura \ref{fig:vista_distrib}). \\
        \hline
            \textbf{Justificación} & \pendiente{completar la justificación.} \\
        \hline
            \textbf{Decisiones Relacionadas} & DA01 \\
        \hline
    \end{tabular}
\end{table}

\begin{table}[H]
    \small
    \centering
    \caption{DA06: Componente móvil nativo para Android}
    \label{tab:da06}
    \begin{tabular}{|p{4.1cm}|p{11.7cm}|}
        \hline
            \textbf{Identificador} & DA06 \\
        \hline
            \textbf{Nombre} & Componente móvil nativo para Android \\
        \hline
            \textbf{Categorías} & Componente móvil; AC18, AC19, AC20. \\
        \hline
            \textbf{Problema} & El chofer debe poder consultar la Guía y registrar los eventos del viaje sin conectividad, y la sincronización posterior debe ser idempotente. \\
        \hline
            \textbf{Alternativas} & \pendiente{completar las alternativas consideradas.} \\
        \hline
            \textbf{Decisión} & Aplicación nativa para Android en Kotlin con Jetpack Compose, persistencia local con Room y un identificador UUID generado en el dispositivo para cada evento pendiente. La sincronización corre en segundo plano con WorkManager, que espera a tener conexión y reintenta con esperas crecientes, y las llamadas a la API REST usan Retrofit con OkHttp. \\
        \hline
            \textbf{Justificación} & \pendiente{completar la justificación.} \\
        \hline
            \textbf{Decisiones Relacionadas} & DA04 \\
        \hline
    \end{tabular}
\end{table}

\begin{table}[H]
    \small
    \centering
    \caption{DA07: Posiciones en MongoDB con idempotencia por índice único}
    \label{tab:da07}
    \begin{tabular}{|p{4.1cm}|p{11.7cm}|}
        \hline
            \textbf{Identificador} & DA07 \\
        \hline
            \textbf{Nombre} & Posiciones en MongoDB con idempotencia por índice único \\
        \hline
            \textbf{Categorías} & Ingesta de eventos; opcional 10; R6; AC08, AC15, AC16. \\
        \hline
            \textbf{Problema} & Las posiciones son el dato de mayor volumen y solo se agregan. Además, un mismo evento de tracking puede recibirse más de una vez (reentregas del broker o reenvíos del nodo) y la cola tiene varios consumidores en paralelo. Un evento duplicado no debe alterar el estado del viaje ni generar casos adicionales. \\
        \hline
            \textbf{Alternativas} & Guardar las posiciones en PostgreSQL con una restricción de unicidad; usar una colección de series temporales de MongoDB. \\
        \hline
            \textbf{Decisión} & Guardar las posiciones en una colección normal de MongoDB, con un índice único sobre el par (nodo, identificador del evento) e índices por viaje y por matrícula con el momento de generación. Un evento que viola el índice único se descarta como duplicado. Como MongoDB no participa de la transacción del mensaje, cada posición se guarda sin procesar y se marca como procesada cuando se confirma esa transacción. Si llega de nuevo una posición que quedó sin procesar, porque un intento anterior falló después de guardarla, no se guarda otra vez pero se completan los pasos que siguen al guardado; si ya estaba procesada, se descarta. Las posiciones referencian el vehículo, el viaje y el nodo por su identificador, sin claves foráneas. \\
        \hline
            \textbf{Justificación} & Las posiciones son el dato que más crece en la plataforma, y guardarlas aparte evita que degraden a PostgreSQL. Se usa una colección normal de MongoDB y no una de series temporales, porque estas no admiten índices únicos para descartar duplicados ni permiten modificar una posición después para asociarla a su viaje. Si una posición llega de nuevo, el índice único la identifica como duplicada y no se guarda dos veces. \\
        \hline
            \textbf{Decisiones Relacionadas} & DA02, DA08 \\
        \hline
    \end{tabular}
\end{table}

\begin{table}[H]
    \small
    \centering
    \caption{DA08: Registro de nodos periféricos con estados}
    \label{tab:da08}
    \begin{tabular}{|p{4.1cm}|p{11.7cm}|}
        \hline
            \textbf{Identificador} & DA08 \\
        \hline
            \textbf{Nombre} & Registro de nodos periféricos con estados \\
        \hline
            \textbf{Categorías} & Servicios de negocio; CU-19; CU-10. \\
        \hline
            \textbf{Problema} & Los administradores deben poder gestionar los nodos periféricos (CU-19), y la ingesta de tracking solo debe aceptar eventos de un nodo registrado y habilitado, que es la precondición del CU-10. \\
        \hline
            \textbf{Alternativas} & \pendiente{completar las alternativas consideradas.} \\
        \hline
            \textbf{Decisión} & Un registro de nodos con tipo (balanzas o tracking), punto de acceso y estado (deshabilitado, habilitado o baja). Todo nodo nuevo arranca deshabilitado, cambiar el punto de acceso lo vuelve a deshabilitar y la baja es lógica. \\
        \hline
            \textbf{Justificación} & \pendiente{completar la justificación.} \\
        \hline
            \textbf{Decisiones Relacionadas} & DA07 \\
        \hline
    \end{tabular}
\end{table}

\begin{table}[H]
    \small
    \centering
    \caption{DA09: Autorización declarativa en la capa de servicios}
    \label{tab:da09}
    \begin{tabular}{|p{4.1cm}|p{11.7cm}|}
        \hline
            \textbf{Identificador} & DA09 \\
        \hline
            \textbf{Nombre} & Autorización declarativa en la capa de servicios \\
        \hline
            \textbf{Categorías} & Seguridad; servicios de negocio; R9; AC03. \\
        \hline
            \textbf{Problema} & La autorización debe aplicarse en la capa de servicios y no solo en la interfaz de usuario (R9). Además, un responsable o un chofer solo debe acceder a los datos de su propia empresa. \\
        \hline
            \textbf{Alternativas} & \pendiente{completar las alternativas consideradas.} \\
        \hline
            \textbf{Decisión} & Cada request se autentica con un mecanismo de Jakarta Security que valida el JWT de la cookie o del encabezado \texttt{Authorization} y le informa al contenedor el usuario y sus roles. Los EJB declaran los roles permitidos con \texttt{@RolesAllowed} y verifican además que el usuario tenga acceso a la empresa sobre la que opera. \\
        \hline
            \textbf{Justificación} & \pendiente{completar la justificación.} \\
        \hline
            \textbf{Decisiones Relacionadas} & DA03 \\
        \hline
    \end{tabular}
\end{table}

\begin{table}[H]
    \small
    \centering
    \caption{DA10: Consulta a balanzas al finalizar el viaje}
    \label{tab:da10}
    \begin{tabular}{|p{4.1cm}|p{11.7cm}|}
        \hline
            \textbf{Identificador} & DA10 \\
        \hline
            \textbf{Nombre} & Consulta a balanzas al finalizar el viaje \\
        \hline
            \textbf{Categorías} & Integración; fiscalización automática; R3. \\
        \hline
            \textbf{Problema} & La regla de exceso de carga necesita los pesajes del vehículo, pero el nodo de balanzas no los envía por iniciativa propia: hay que definir en qué momento se le consultan (CU-11). \\
        \hline
            \textbf{Alternativas} & Consultar con cada evento de posición; consultar periódicamente durante el viaje; que el nodo de balanzas notifique cada pesaje; consultar al finalizar el viaje. \\
        \hline
            \textbf{Decisión} & La integración con el Sistema de Balanzas es una consulta iniciada por el componente central, como exige la letra. El Motor de Fiscalización consulta los pesajes cuando procesa el fin de un viaje, por el vehículo y cada fecha del viaje. \\
        \hline
            \textbf{Justificación} & Se descartó que el nodo de balanzas notificara cada pesaje porque la letra define la integración como solicitud--respuesta (R3): el componente central pide los pesajes de un vehículo en una fecha y el nodo devuelve todas las pesadas de ese vehículo en esa fecha. Consultar con cada posición o periódicamente repetiría la misma consulta muchas veces, porque cada respuesta trae las pesadas de todo el día. Al finalizar el viaje ya se conocen el vehículo y las fechas, y alcanza con una consulta por fecha. Si Balanzas no responde, la evaluación queda pendiente sin detener la plataforma. \\
        \hline
            \textbf{Decisiones Relacionadas} & DA01, DA08 \\
        \hline
    \end{tabular}
\end{table}

\begin{table}[H]
    \small
    \centering
    \caption{DA11: Permisos aprobados por un funcionario}
    \label{tab:da11}
    \begin{tabular}{|p{4.1cm}|p{11.7cm}|}
        \hline
            \textbf{Identificador} & DA11 \\
        \hline
            \textbf{Nombre} & Permisos aprobados por un funcionario \\
        \hline
            \textbf{Categorías} & Servicios de negocio; CU-06; CU-07; CU-14. \\
        \hline
            \textbf{Problema} & La letra pide que el responsable cargue el Permiso Nacional de Circulación con su número y período de validez, y que los funcionarios fiscalicen los permisos. Si el permiso habilitara al vehículo apenas se carga, el responsable podría declarar cualquier fecha de fin y el vehículo quedaría habilitado para asignarle viajes sin ningún control. \\
        \hline
            \textbf{Alternativas} & Que el permiso habilite apenas se carga y se controle después en la fiscalización; que el responsable cargue todo el período y el funcionario solo apruebe o rechace; que el responsable cargue el número y la fecha de inicio y el funcionario, al aprobar, fije la fecha de fin. \\
        \hline
            \textbf{Decisión} & El permiso tiene estado: pendiente, aprobado o rechazado. El responsable carga el número y la fecha de inicio y el permiso queda pendiente, sin fecha de fin. Un funcionario lo revisa: lo aprueba fijando la fecha de fin, y puede corregir la de inicio, o lo rechaza con un motivo. Solo un permiso aprobado y vigente habilita al vehículo. Un permiso aprobado no se modifica ni se elimina; para renovarlo se carga uno nuevo. Un rechazado se puede corregir y volver a enviar. Las operaciones del responsable reciben solo los datos que puede declarar, así que la fecha de fin y el estado no se pueden fijar desde la API aunque se envíen. \\
        \hline
            \textbf{Justificación} & En la realidad el número y el vencimiento del permiso los fija el MTOP, no la empresa; la empresa solo informa un permiso ya emitido. Dejar la fecha de fin a cargo del funcionario evita que el responsable la declare a su conveniencia, y la aprobación deja registrado quién revisó cada permiso y cuándo. Se descartó controlarlo solo en la fiscalización porque el vehículo ya habría viajado. El estado sigue el mismo criterio que los nodos periféricos (DA08), que arrancan deshabilitados hasta que un administrador los habilita. \\
        \hline
            \textbf{Decisiones Relacionadas} & DA08, DA09 \\
        \hline
    \end{tabular}
\end{table}

\subsection{Decisiones Pendientes}

La Tabla \ref{tab:decisiones-pendientes} lista lo que la letra pide definir y justificar, junto con otras decisiones que surgieron durante el desarrollo, que todavía no están definidas.

\begin{longtable}{|p{4.6cm}|p{7.4cm}|p{3.5cm}|}
    \caption{Decisiones pendientes}\label{tab:decisiones-pendientes}\\
    \hline
    \textbf{Tema} & \textbf{Qué se necesita definir} & \textbf{Dónde impacta} \\
    \hline
    \endfirsthead
    \hline
    \textbf{Tema} & \textbf{Qué se necesita definir} & \textbf{Dónde impacta} \\
    \hline
    \endhead
    Volumen de tracking & Cantidad de vehículos reportando en forma concurrente, frecuencia de reporte y picos de al menos tres veces ese volumen. & AC08, prueba de performance, generador de carga. \\
    \hline
    Reintentos y cola de mensajes muertos & Espera entre reintentos, cantidad de intentos y qué se hace con los mensajes que terminan en la cola de mensajes muertos. & DA02, AC12. \\
    \hline
    Servicio de rutas y geocodificación & Qué servicio externo se usa para ubicar el origen y el destino y obtener el recorrido esperado, con sus límites de uso. & CU-07, Vista Lógica, Distribución. \\
    \hline
    Criterio de desvío de ruta & Distancia máxima al recorrido esperado, tiempo fuera de él y cantidad mínima de posiciones, con sus valores por defecto. & CU-11, parametrización de reglas. \\
    \hline
    Momento de evaluación del desvío & Si el desvío se evalúa al finalizar el viaje o con cada posición recibida. & CU-11, Motor de Fiscalización. \\
    \hline
    Umbral de agregación & Umbral mínimo de agregación de la información pública. & AC06, R12. \\
    \hline
    Escalado y orquestador & Dónde corre el orquestador del opcional 4 y cómo escala con varias réplicas la sesión del backoffice JSF. & AC07, R14, DA02, Vista de Distribución. \\
    \hline
    Tecnología del backoffice & Si el backoffice queda en JSF o pasa a React. & DA05, AC07. \\
    \hline
    Plataforma de los periféricos & Qué PaaS o IaaS ejecuta los nodos periféricos y el suplantador de la PDI. & R15, Vista de Distribución. \\
    \hline
    Plataforma de observabilidad & En qué PaaS o IaaS se ejecuta la Plataforma de Observabilidad en el despliegue. & AC21 a AC24, Vista de Distribución. \\
    \hline
    Fotografías de los incidentes & Dónde se guardan y cómo se envían desde el móvil al recuperar la conexión. & CU-09, Persistencia. \\
    \hline
    Plazo de descargo & Si el plazo se cuenta en días hábiles o corridos. & CU-11, parametrización de reglas. \\
    \hline
    Viaje sin fin registrado & Cómo se evalúa el exceso de carga si el chofer nunca registra el fin del viaje, y cómo se cubre la regla de viaje sin Guía por pesaje. & DA10, CU-11. \\
    \hline
    Prueba de performance & Herramienta y entorno en el que se ejecuta. & AC09 a AC11. \\
    \hline
\end{longtable}
